\documentclass[12pt]{article}
\usepackage{graphicx} % Required for inserting images
\usepackage[margin=0.5in]{geometry}
\usepackage{amsmath, amssymb}
\usepackage{mathrsfs}


\usepackage{soul}


% Dr. David Brown Commands
\newcommand{\ds}{\displaystyle}
\newcommand{\R}{\mathbb{R}}
\newcommand{\N}{\mathbb{N}}
\newcommand{\Q}{\mathbb{Q}}
\newcommand{\C}{\mathbb{C}}


% Commands to get script r
\usepackage{calligra}
\DeclareMathAlphabet{\mathcalligra}{T1}{calligra}{m}{n}
\DeclareFontShape{T1}{calligra}{m}{n}{<->s*[2.2]callig15}{}
\newcommand{\scripty}[1]{\ensuremath{\mathcalligra{#1}}}

% Unit commands
\newcommand{\degree}{\text{°}}
\newcommand{\unitSpeed}{\frac{\text{m}}{\text{s}}}
\newcommand{\unitAcceleration}{\frac{\text{m}}{\text{s}^2}}

% Constant commands
\newcommand{\avogadro}{6.022\times10^{23}}

\newcommand{\boltzmannConstK}{1.381\times10^{-23}\frac{\text{J}}{\text{K}}}
\newcommand{\planckConst}{6.626\times10^{-34}\text{J}\cdot\text{s}}

\newcommand{\atmP}{1.01325\times10^5\,\text{Pa}}

% Commands to easily write partial derivatives
\newcommand{\partDeriv}[2]{\frac{\partial #1}{\partial #2}}
\newcommand{\doublePartDeriv}[2]{\frac{\partial^2 #1}{\partial^2 #2}}


\title{Problem 85}
\author{Gabriel Decker}


\begin{document}



\maketitle
\begin{flushleft}

For this problem, we can consider just the 2nd half of a vertical jump, so only in the y direction.
We consider only the force of gravity on the person jumping.
This problem being a constant-force problem where someone would jump with some initial speed implies that the trajectory of the jump is the same going up as it is going down.
Therefore, we can simplify our problem by just considering the top of the jump down to the floor.
So consider $t_0=0$ as the time when the jumper is at rest in the vertical direction.
It's at the height $y=h$.\\~\\

We'll use the following equations.
$$v=v_0+at$$
$$y=y_0+v_0t+\frac{1}{2}at^2$$\\~\\

At $t=0$, the top of the jump, $v_0=0$.
The system is contant-acceleration under gravity, so $a=-g$.
Therefore, the above equations become the following.

$$v=-gt$$
$$y=h-\frac{1}{2}gt^2$$\\~\\

With this setup, we can find the time when the person hits the ground and when the person is halfway down ($y=\frac{1}{2}h)$) by plugging in those $y$ values and solving for $t$.
First, $y=\frac{1}{2}h$.
$$\frac{1}{2}h=h-\frac{1}{2}gt^2$$
$$\implies \frac{1}{2}gt^2=\frac{1}{2}h$$
$$\implies t^2=\frac{h}{g}$$
$$\implies t=\sqrt{\frac{h}{g}}$$\\~\\

Then, $y=0$.
$$0=h-\frac{1}{2}gt^2$$
$$\implies \frac{1}{2}gt^2=h$$
$$\implies t^2=\frac{2h}{g}$$
$$\implies t=\sqrt{\frac{2h}{g}}$$\\~\\

The time for $y=0$ is the total time of the descent.
We can determine the fraction of the time spent in the upper half of the trajectory by dividing the second from the first.
$$\frac{\sqrt{h/g}}{\sqrt{2}\sqrt{h/g}}=\frac{1}{\sqrt{2}}\approx0.707=70.7\%$$

\end{flushleft}

\end{document}
